\documentclass{article}
\usepackage{amsmath}
\usepackage{amssymb}
\usepackage{physics}
\usepackage{geometry}
\geometry{margin=1in}

\title{Symbolic Orbital Transfer System}
\author{Albert et al.}
\date{\today}

\begin{document}

\maketitle

\section*{Overview}
This document describes the canonical symbolic system for orbital transfers constructed by Albert, generalizing classical two-body mechanics to an arbitrary multi-attractor environment, with spline-based thrust formulation over arc length.

It provides a complete symbolic family of constraints suitable for subsequent numerical optimization, machine learning control, or analytical exploration.

\section*{Equations of Motion}
Let $r(s)$ denote the position vector parameterized by arc length $s$, and $F_\text{extra}(s)$ the applied spline force. 

With multiple gravitational centers $\{(c_i, \mu_i)\}$ where $c_i$ is the center position and $\mu_i$ is the gravitational parameter, the force contributions are
\[
F_{\text{grav},i} = -\mu_i \frac{r - c_i}{\norm{r - c_i}^3}
\]

Then the symbolic arc-length equation of motion is
\[
\frac{d^2 r}{ds^2} = \sum_i F_{\text{grav},i} + F_\text{extra}(s)
\]

\section*{Total Mechanical Energy}
The symbolic total energy along the trajectory is given by
\[
E_{\text{total}}(s) = \frac12 \left\|\frac{dr}{ds}\right\|^2 - \sum_i \frac{\mu_i}{\norm{r - c_i}}
\]

\section*{Integral Force Cost}
The total symbolic cost over the distance traveled is represented by
\[
\int_0^L \norm{F_\text{extra}(s)} \, ds
\]
where $L$ is the total arc length of the trajectory.

\section*{Boundary Conditions}
We enforce the symbolic initial and terminal constraints
\[
r(0) = r_{\text{start}}, \quad r(L) = r_{\text{end}}
\]

\section*{Remarks}
\begin{itemize}
    \item All quantities are purely symbolic, defined via the \texttt{sympy} library.
    \item The system does not solve for a trajectory explicitly; it produces the symbolic family of solutions subject to these constraints.
    \item This allows deferred optimization, variational calculus, or machine learning exploration on the exact same symbolic landscape.
\end{itemize}

\end{document}
